\documentclass[10pt,a4paper]{amsart}
\usepackage[margin=19mm]{geometry}
\usepackage{amsmath,amssymb,mathtools}
\usepackage[scr=rsfs,cal=euler]{mathalfa}
\usepackage{xfrac}
\usepackage[unicode,hidelinks,bookmarks=false]{hyperref}
\newcommand{\ie}{i.e.\ }
\newcommand{\cf}{cf.\ }
\newcommand{\resp}{respectively\ }
\newcommand{\Subsection}{\subsection}
\providecommand{\nobreakdash}{\nobreak\hbox{-}}
\setlength{\emergencystretch}{2em}
\multlinegap0pt
\usepackage{bm}
\usepackage{bbm}
\usepackage{dsfont}
\usepackage{xspace}
\usepackage{cancel}
\usepackage{mathtools}
\usepackage{amsthm}
\makeatletter
\@ifundefined{theorem}{\newtheorem{theorem}{Theorem}}{}
\@ifundefined{lemma}{\newtheorem{lemma}[theorem]{Lemma}}{}
\makeatother
\newcommand{\Ft}{\mathbb{F}_2}
\newcommand{\spann}[1]{\langle #1\rangle}
\title[Non-isomorphic restricted Lie algebras with isomorphic restr]{Non-isomorphic restricted Lie algebras with isomorphic restricted enveloping algebras}
\author{Xabier Garc\'ia-Mart\'inez}
\date{Source: arXiv:2606.17282v1 (CC BY 4.0)}
\begin{document}
\maketitle

\noindent\emph{Source and licence.} Non-isomorphic restricted Lie algebras with isomorphic restricted enveloping algebras. Version arXiv:2606.17282v1 (CC BY 4.0), distributed under the Creative Commons Attribution 4.0 International licence, as recorded in the arXiv licence metadata for this exact version and shown on the abstract page https://arxiv.org/abs/2606.17282v1. Full rights notice: licenses/arxiv-2606.17282-CC-BY-4.0.txt.\par\smallskip

\noindent\textit{Editorial scope.} Counted unit: the Theorem whose source label is \texttt{thm:char2} in the article (source0.tex lines 368--377, source label \texttt{thm:char2}), delivered together with the article's own complete original proof of it (source0.tex lines 379--470, one \texttt{proof} environment of 92 lines). No mathematics is written, repaired or re-derived: statement and proof are the article's own text line for line. Required context retained uncounted: eq:square (lines 172--176); eq:pres (lines 202--206); lem:valid2 (lines 252--260). Every definition, display and label of those lines is retained unchanged. Omitted from this package and not redistributed: 1--171, 177--201, 207--251, 261--367, 378--378, 471--901. Nothing inside a retained excerpt is omitted or altered: every retained body is delivered exactly as the article's lines are written, so no omission receipt of any kind accompanies this package. Citations: the retained excerpts carry no citation command at all, so no bibliography entry of the article is transcribed and no reference is redistributed. Reference adaptation: none was needed and none was made; every reference inside the delivered unit resolves inside it, so no unresolved reference or citation is produced. Printed slips kept literal: no printed word, symbol, hypothesis or number of the counted statement or of its proof was altered. Layout and class: the article's own class is \texttt{amsart} with packages including amsmath, amssymb, amsthm, mathtools, geometry, xcolor, hyperref, enumitem; the wrapper is the lane's pinned \texttt{amsart} editorial wrapper at 10pt on A4, which restates the theorem-like environments the retained text needs and adds the article's own macro definitions that the retained text needs (\texttt{\textbackslash Ft}, \texttt{\textbackslash spann}), with extra wrapper packages (bm, bbm, dsfont, xspace, cancel, mathtools). The delivered numbering is the unit's own. Assets: no retained span contains a figure environment, a graphics-inclusion command, a PDF-page inclusion, a table or a TikZ picture; no image file is copied into this package, the package declares no asset dependency and no image file is redistributed with it. Licence: version arXiv:2606.17282v1 of this article is distributed under the Creative Commons Attribution 4.0 International licence, as recorded in the arXiv licence metadata for this exact version and shown on the abstract page https://arxiv.org/abs/2606.17282v1. A full rights notice is delivered at licenses/arxiv-2606.17282-CC-BY-4.0.txt. Characters: every retained excerpt is pure ASCII, so the delivered engine, XeLaTeX with its default Latin Modern text font, reports no missing character.
\begin{equation}\label{eq:square}
(r+r')^p=r^p+r'^p \text{ whenever } [r,r']=0,
\text{ and, for } p=2,
\Bigl(\sum_i r_i\Bigr)^{2}=\sum_i r_i^{2}+\sum_{i<j}\,[r_i,r_j].
\end{equation}

\begin{equation*}
u(\mathfrak{g})\;=\;F\langle e_1,\dots,e_n\rangle\Big/
\bigl(\,e_ie_j-e_je_i-[e_i,e_j],\quad e_i^{\,p}-e_i^{[p]}\,\bigr)_{i,j}\,.
\tag{P}\label{eq:pres}
\end{equation*}

\begin{lemma}\label{lem:valid2}
$L_k$ and $H_k$ are restricted Lie algebras of dimension $n=k+4$, generated
by $a,b$ and by $x,y$ respectively. Both have nilpotency class $3$, with
lower central dimension sequences $(n,2,1,0)$ for $L_k$ and $(n,3,2,0)$ for
$H_k$. In particular
\[
L_k'=\spann{v,b_4},\qquad H_k'=\spann{w,y_4,c}.
\]
\end{lemma}

\begin{theorem}\label{thm:char2}
For every $k\ge3$ the assignment
\[
\varphi(a)=X\coloneqq x,\qquad
\varphi(b)=Y\coloneqq x+y+x\,g
= x+y+x\bigl(y^{2^{k}-4}+y^{2^{k+1}-8}\bigr),
\]
extends to an isomorphism of $\Ft$-algebras
$\varphi\colon u(L_k)\to u(H_k)$.
\end{theorem}

\begin{proof}
	We define the images of the basis vectors of $L_k$ as the following words in $X$ and $Y$:
\begin{gather*}
f_a\coloneqq X,\qquad f_b\coloneqq Y,\qquad f_{b_2}\coloneqq Y^2,\qquad f_{a_2}\coloneqq X^2,\qquad
f_v\coloneqq [X,Y],\\
f_{b_{2^j}}\coloneqq Y^{2^j}\ (2\le j\le k-1),\qquad f_d\coloneqq X^4 .
\end{gather*}
The assignment extends to an
algebra homomorphism $\varphi\colon u(L_k)\to u(H_k)$ with $\varphi(e)=f_e$
for every basis vector $e$ of $L_k$ as soon as the elements $f_e$ satisfy all
the relations \eqref{eq:pres} of $u(L_k)$. Let us verify them:

Let us begin by checking the power relations. By construction $f_a^{\,2}=f_{a_2}$,
$f_b^{\,2}=f_{b_2}$, $f_{b_2}^{\,2}=f_{b_4}$,
$f_{b_{2^j}}^{\,2}=f_{b_{2^{j+1}}}$ for $2\le j\le k-2$,
$f_{a_2}^{\,2}=f_d$ and $f_d^{\,2}=X^8=x^8=0$, matching the $2$-map of~$L_k$. There are two remaining power relations: $f_v^{\,2}=0$ and
$f_{b_q}^{\,2}=f_d$.

Since $g$ is central, $[x,xg]=g\,[x,x]=0$, so
\[
f_v=[X,Y]=[x,y]+[x,xg]=w,
\qquad\text{whence}\qquad
f_v^{\,2}=w^2=0.
\]

By \eqref{eq:square} applied to $Y=x+y+xg$: the
squares of the summands are $x^2$, $y^2$ and $(xg)^2=x^2g^2$, and the
cross-brackets are $[x,y]=w$, $[x,xg]=0$ and $[y,xg]=g\,[y,x]=gw$. Hence
\begin{equation}\label{eq:Ysq}
Y^2=x^2+y^2+x^2g^2+(1+g)\,w .
\end{equation}

The four summands of \eqref{eq:Ysq} commute
pairwise: $g$ is central and $[x^2,y^2]=[x^2,w]=[y^2,w]=0$. So by \eqref{eq:square} their squares add up:
\[
Y^4=x^4+y^4+x^4g^4+(1+g)^2w^2=\delta+x^4g^4=\delta.
\]
Thus $f_{b_4}=\delta$ and
$Y^8=\delta^2=y^8$, so that
\[
f_{b_{2^j}}=Y^{2^j}=\bigl(Y^8\bigr)^{2^{j-3}}=y^{2^j}\quad(3\le j\le k-1):
\]
all the elements $f_{b_4},\dots,f_{b_q},f_d=x^4$ are central in $u(H_k)$. The
missing power relation follows:
\[
f_{b_q}^{\,2}=Y^{2^k}=\bigl(Y^{8}\bigr)^{2^{k-3}}
=\bigl(y^{8}\bigr)^{2^{k-3}}=y^{2^k}=x^4=f_d.
\]

Let us proceed now to check the brackets. In $L_k$ every basis vector other than
$a,b,b_2,a_2,v$ is central, and the corresponding images
$f_{b_4},\dots,f_{b_q},f_d$ are central in $u(H_k)$, hence all relations~\eqref{eq:pres} involving them hold, both sides being zero. Moreover
$[f_a,f_b]=f_v$ holds by construction, and
$[f_a,f_{a_2}]=[x,x^2]=0$, $[f_b,f_{b_2}]=[Y,Y^2]=0$ trivially. Let us compute the remaining pairs:
\begin{align*}
[f_a,f_{b_2}]&=[x,Y^2]=[x,y^2]+g^2\,[x,x^2]+(1+g)\,[x,w]
=y^4+(1+g)\,\delta\\
&=(y^4+\delta)+g\delta=x^4+x^4=0=f_{[a,b_2]},\\[2pt]
[f_{a_2},f_b]&=[x^2,Y]=[x^2,x]+[x^2,y]+g\,[x^2,x]=\delta
=f_{b_4}=f_{[a_2,b]},\\[2pt]
[f_a,f_v]&=[x,w]=\delta=f_{b_4}=f_{[a,v]},\\[2pt]
[f_b,f_v]&=[Y,w]=[x,w]+[y,w]+g\,[x,w]=\delta+y^4+g\delta
=\delta+y^4+x^4=\delta+\delta=0=f_{[b,v]},\\[2pt]
[f_{b_2},f_{a_2}]&=[Y^2,x^2]=0,\qquad
[f_{b_2},f_v]=[Y^2,w]=0,\qquad
[f_{a_2},f_v]=[x^2,w]=0.
\end{align*} 
In each case the value agrees with
the image of the corresponding bracket of $L_k$. All relations~\eqref{eq:pres} of $u(L_k)$
hold, so $\varphi$ is a well-defined algebra homomorphism.

To see the bijectivity, by Lemma~\ref{lem:valid2} and the
Poincar\'e--Birkhoff--Witt theorem the two algebras have dimension $2^{\,n}$,
so it suffices to show that $\varphi$ is surjective. As $L_k$ is generated by~$a,b$, 
so is $u(L_k)$ as an associative algebra, hence the image of $\varphi$
is the subalgebra $B$ generated by~$X=x$ and $Y$, and it is enough to show
$x,y\in B$. Certainly $x=X\in B$, and, the characteristic being two,
\[
\eta\coloneqq X+Y=y+xg\in B .
\]
Moreover, $\eta^4=y^4$. Indeed $X^2=x^2$ and $[X,Y]=w$, so \eqref{eq:square} gives $\eta^2=X^2+Y^2+[X,Y]=x^2+Y^2+w$. Squaring
once more, the cross terms vanish, since $w^2=0$,
$[x^2,w]=0$, and $[x^2,Y^2]=[Y^2,w]=0$ were shown
among the bracket relations above, whence
\[
\eta^4=X^4+Y^4=x^4+\delta=y^4 ,
\]
using $Y^4=\delta=x^4+y^4$. Now $g=(y^4)^{m}+(y^4)^{2m}$ with
$m\coloneqq 2^{k-2}-1$ (because $2q-4=4m$ and $4q-8=8m$), so
$g=(\eta^4)^{m}+(\eta^4)^{2m}\in B$. Therefore $xg\in B$ and
$y=\eta+xg\in B$. Thus $B=u(H_k)$, and $\varphi$ is an isomorphism.
\end{proof}

\end{document}
